Semántica no reflejada:
\begin{enumerate}
    \item \textbf{Cuando se genera un movimiento de venta con evento, el evento se guardará exclusivamente en la tabla de ventas, y no en la tabla de eventos.}
    \item Cantidad de movimiento no puede ser negativa ni superior a la cantidad de stock modificado.
\end{enumerate}

\begin{lstlisting}[caption={Modelo Relacional - Papertory}, label={fig:modelo-relacional-papertory}]

    USUARIO (id_usuario, nombre, email, password, Rol)
        CP: id_usuario

    PROVEEDOR (id_proveedor, nombre, email, vatId, telefono, direccion, ciudad, pais)
        CP: id_proveedor

    LOCALIZACION (id_localizacion, nombre)
        CP: id_localizacion

    PRODUCTO (id_producto, nombre, descripcion, codigo, activo, tipo)
        CP: id_producto

    MATERIALPROMOCIONAL (id_producto, tipo, ancho, alto, descripcion)
        CP: id_producto
        CE: id_producto REFERENCIA PRODUCTO(id_producto)

    MERCHANDISING (id_producto, talla, precio, iva, descripcion)
        CP: id_producto
        CE: id_producto REFERENCIA PRODUCTO(id_producto)

    EXISTENCIA (id_existencia, id_producto, id_localizacion, id_proveedor, fecha_entrada, fecha_apertura, cantidad_inicial, cantidad_actual)
        CP: id_existencia
        CE: id_producto REFERENCIA PRODUCTO(id_producto)
        CE: id_localizacion REFERENCIA LOCALIZACION(id_localizacion)
        CE: id_proveedor REFERENCIA PROVEEDOR(id_proveedor)

    EVENTO (id_evento, nombre, descripcion)
        CP: id_evento

    MOVIMIENTO (id_movimiento, id_existencia, id_evento, id_usuario, fecha, cantidad, tipo)
        CP: id_movimiento
        CE: id_existencia REFERENCIA EXISTENCIA(id_existencia)
        CE: id_evento REFERENCIA EVENTO(id_evento)
        CE: id_usuario REFERENCIA USUARIO(id_usuario)

    VENTA (id_venta, id_evento, id_usuario, fecha, metodo)
        CP: id_venta
        CE: id_evento REFERENCIA EVENTO(id_evento)
        CE: id_usuario REFERENCIA USUARIO(id_usuario)

    LINEA (id_linea, id_venta, id_movimiento, iva, precio_unitario, cantidad)
        CP: id_linea, id_venta
        CE: id_venta REFERENCIA VENTA(id_venta)
        CE: id_movimiento REFERENCIA MOVIMIENTO(id_movimiento)



    
\end{lstlisting}


\begin{longtable}{|l|l|p{3.8cm}|p{2.2cm}|c|c|c|c|}
    \hline
    \multicolumn{8}{| c |}{Begin of Table} \\
    \hline
    Entidad o relación & Atributos & Descripción & Tipo de datos & N & M & D & P \\
    \hline
    \multirow{6}{*}{Persona}
    & Nombre & Nombre y apellidos de cada usuario & Carácter variable 100 & NO & NO & NO & NO \\
    \cline{2-8}
    & Id\_usuario \{PK\} & Identificador único de cada usuario & Carácter variable 30 & NO & NO & NO & NO \\
    \cline{2-8}
    & Correo & Dirección de email único de cada usuario & Carácter variable 100 & NO & NO & NO & NO \\
    \cline{2-8}
    & Contraseña & Contraseña de cada usuario para acceder a la aplicación & Carácter variable 20 & NO & NO & NO & NO \\
    \cline{2-8}
    & Rol & Rol de cada usuario (administrador, empleado, cliente) & Carácter variable 20 & NO & NO & NO & NO \\
    \cline{2-8}
    & tipo\_usuario & Tipo de usuario: administrador, usuario & Carácter variable 20 & NO & NO & NO & 'Normal' \\
    \hline
    \hline
    \caption{This is the caption for the test table.}
\end{longtable}\label{tab:diccionario-datos}
